% ---------------------------------------------------------------------
%  Estructura: partes, problemas, incisos, pasos
% ---------------------------------------------------------------------

\makeatletter
% \parte{I}{Decaimiento radiactivo}{20}
% Cada parte y cada problema empiezan en hoja nueva; el primer problema de
% una parte se queda en la misma hoja que el título de la parte.
\newif\ifrecienparte
\newcommand{\parte}[3]{%
  \clearpage\phantomsection\global\recienpartetrue
  \section*{Parte #1. #2 (#3 pts)}%
  \addcontentsline{toc}{section}{Parte #1. #2}%
  \markboth{Parte #1. #2}{Parte #1. #2}}

% \problema{3}{Decaimiento con fuente}{6}
% Se puede referir con  Problema~\ref{prob:3}
\newcommand{\problema}[3]{%
  \ifrecienparte\global\recienpartefalse\else\clearpage\fi
  \phantomsection
  \subsection*{Problema #1. #2\hfill[#3 pts]}%
  \addcontentsline{toc}{subsection}{Problema #1. #2}%
  \protected@edef\@currentlabel{#1}%
  \label{prob:#1}}

% \inciso{a}{Esquema CN con fuente}
% Si el título lleva matemáticas, usa \texorpdfstring{$...$}{texto} para el índice del PDF.
\newcommand{\inciso}[2]{%
  \phantomsection
  \subsubsection*{(#1) #2}%
  \addcontentsline{toc}{subsubsection}{(#1) #2}}
\makeatother

% \paso{1}{Generación del mallado}
\newcommand{\paso}[2]{\par\medskip\textbf{Paso #1.} \textit{#2}\par\smallskip}

% ---------------------------------------------------------------------
%  Notas dentro del texto (sin cajas)
% ---------------------------------------------------------------------

% Enunciado del problema
\newenvironment{enunciado}{\par}{\par\medskip}

% porque / ojo / leveque: la explicación va directo como párrafo normal, sin título.
% (El título opcional, p. ej. \begin{porque}[...], se ignora.)
\newenvironment{porque}[1][]{\par}{\par}
\newenvironment{ojo}[1][]{\par}{\par}
\newenvironment{leveque}[1][]{\par}{\par}

% Resultado final destacado: \resultado{ ecuación }
\newcommand{\resultado}[1]{\[\boxed{\displaystyle #1}\]}

% Problema aún sin resolver
\newcommand{\pendiente}{\par\textit{Solución pendiente.}\par}

% ---------------------------------------------------------------------
%  Material generado por Python (no rompe la compilación si falta)
% ---------------------------------------------------------------------
% \figurapython{p3_solucion}{Pie de figura}{fig:p3sol}{0.8}
\newcommand{\figurapython}[4]{%
  \begin{figure}[htbp]\centering
  \IfFileExists{figuras/#1.pdf}{%
    \includegraphics[width=#4\textwidth]{figuras/#1.pdf}}{%
    \fbox{\parbox{0.7\textwidth}{\centering\small
      Falta \texttt{figuras/#1.pdf}.\\ Ejecuta \texttt{make figs} en la raíz.}}}
  \caption{#2}\label{#3}
  \end{figure}}

% \tablapython{p3_d}  — inserta latex/tablas/p3_d.tex
\newcommand{\tablapython}[1]{%
  \IfFileExists{tablas/#1.tex}{\input{tablas/#1.tex}}{%
    \fbox{\small Falta \texttt{tablas/#1.tex}: ejecuta \texttt{make figs}.}}}

% Respuesta de opción múltiple / conceptual: \respuesta{b) porque ...}
% Si se deja vacía, \respuesta{} imprime una línea en blanco.
\newcommand{\respuesta}[1]{%
  \par\smallskip\textbf{Respuesta:}~%
  \if\relax\detokenize{#1}\relax\rule{3cm}{0.4pt}\else #1\fi\par}
